\unnumberedChapter{Введение}

\textbf{Актуальность темы работы}
Работа связана с системами управления вторичными источниками электропитания, их настройкой, отладкой и контролем состояния.

Данная технология предназначена для тестирования и отладки встраиваемого программного обеспечения в части автоматического регулирования координат физического устройства. Автоматическое регулирование координат является распространенной задачей в устройствах реального времени и одной из наиболее ресурсоемкой задач при тестировании и отладке, в том числе требующей высокую квалификацию специалиста. В связи с этим, разрабатываемая в данном проекте технология крайне значима на практике. Технология не только позволяет отказаться от внешних измерительных устройств — анализаторов цепей (Network Analyzers), но и выполнять настройку автоматически, т.е. снизив требования к квалификации специалиста и снизив тем самым стоимость операции.


В системах управления физическими объектами с линейной системой автоматического регулирования наиболее полную и полезную с практической точки зрения информацию о динамических свойствах системы содержат в себе ее частотные характеристики, измеренные экспериментальным путем. В связи с этим, процедуры отладки, настройки и контроля качества микропроцессорных устройств управления реального времени (real-time embedded control systems) зачастую сопровождаются измерением частотных характеристик разомкнутого и замкнутого контура регулирования, объекта управления, регулятора, канала обратной связи и других звеньев системы автоматического регулирования.

В системах управления вторичных источников электропитания, источников бесперебойного питания, регулируемых электроприводов, автономных систем генерирования электроэнергии и многих других систем силовой электроники измерения частотных характеристик проводятся специализированным лабораторным оборудованием, подключаемым в разрыв штатных измерительных цепей системы управления. Потребность в лабораторном оборудовании и в его подключении в разрыв цепей обратной связи делают такой способ измерения непрактичным для систем, находящихся в промышленной эксплуатации.

Цель данной работы заключается в создании алгоритма внутрисхемного измерения частотных характеристик электронных устройств управления физическими объектами, не требующего подключения внешнего оборудования. Возможность физической реализации подобного способа измерения обусловлена специфическими особенностями рассматриваемых в данной работе систем силовой электроники. В этих системах возможно задействовать канал управляющих воздействий как точку ввода программно генерируемого возмущения. При этом физические переменные, характеризующие динамическое состояние системы, доступны к считыванию в реальном времени из периферийных модулей аналого-цифрового преобразования сигналов устройства управления.


\textbf{Цели и задачи работы}

\textbf{Цель работы} -- разработать алгоритмы внутрисхемного измерения частотных характеристик и идентификации параметров вторичных источников электропитания.

\textbf{Задачи работы:}
\begin{enumerate}
    \item Разработать алгоритмы внутрисхемного измерения частотных характеристик.
    \item Разработать алгоритмы идентификации параметров на основе измеренных частотных характеристик.
    \item Выполнить экспериментальную оценку показателей качества измерений и идентификации параметров.
\end{enumerate}

\textbf{Положения, выносимые на защиту}

\textbf{Основные результаты работы}
\begin{enumerate}
    \item Алгоритм внутрисхемного измерения, основанный на, достигающий.
    \item Алгоритм идентификации параметров, основанный на архитектуре автоэнкодера.
\end{enumerate}

\textbf{Достоверность и обоснованность полученных результатов}
Результаты проверены экспериментально на источнике питания.

\textbf{Практическая значимость результатов работы}
Программа «Digital Points» будет применяться в электронной промышленности на этапах тестирования и отладки разрабатываемого встраиваемого программного обеспечения. Области непосредственного применения: производство электронного и электрического оборудования, разработка аппаратно-программных комплексов. В части функционального назначения предлагаемый программный продукт предназначен к применению во время отладочных испытаний при конструировании новых образцов электронных устройств или модернизации серийно выпускаемых изделий преимущественно в качестве автоматизированного средства настройки. На этапе серийного производства отлаженных изделий предлагаемый программный продукт найдет применение для автоматизированной калибровки изделий и контроля качества.

\textbf{Соответствие паспорту специальности}

\textbf{Апробация результатов работы}

\begin{enumerate}
    \item Свидетельство о гос. регистрации программы для ЭВМ. Программный модуль измерения частотных характеристик / \textbf{Р.Л.~Горбунов}, Н.А.~Севостьянов, А.Б.~Дамдинов, И.В.~Александров. –– №~2021665973 от 06.10.2021 (Российская Федерация).
    \item Свидетельство о гос. регистрации программы для ЭВМ. Модуль определения частоты при внутрисхемном измерении частотных характеристик / Н.А.~Севостьянов, \textbf{Р.Л.~Горбунов}. –– №~2024613044 от 08.02.2024 (Российская Федерация).
    \item Свидетельство о гос. регистрации программы для ЭВМ. Модуль построения графиков частотных характеристик / \textbf{Р.Л.~Горбунов}. –– №~2024613165 от 08.02.2024 (Российская Федерация).
    \item Свидетельство о гос. регистрации программы для ЭВМ. Модуль инициализации параметров измерения частотных характеристик / \textbf{Р.Л.~Горбунов}. –– №~2023687756 от 18.12.2023 (Российская Федерация).
    \item Свидетельство о гос. регистрации программы для ЭВМ. Модуль преобразования временных рядов в частотные характеристики / \textbf{Р.Л.~Горбунов}. –– №~2023688842 от 25.12.2023 (Российская Федерация).
    \item Свидетельство о гос. регистрации программы для ЭВМ. Модуль преобразования комплексных переменных в частотные характеристики / \textbf{Р.Л.~Горбунов}. –– №~2023685884 от 30.11.2023 (Российская Федерация).
    \item Свидетельство о гос. регистрации программы для ЭВМ. Классификатор с конфигурируемой архитектурой / \textbf{Р.Л.~Горбунов}. –– №~2026612481 от 28.01.2026 (Российская Федерация).
    \item Свидетельство о гос. регистрации программы для ЭВМ. Функция настройки оптимизатора автоэнкодера / \textbf{Р.Л.~Горбунов}. –– №~2026615254 от 04.02.2026 (Российская Федерация).
\end{enumerate}

\textbf{Личный вклад автора}

\textbf{Объем и структура работы}